% ===================================================================
%  TAULUKKO N:o 5 — Talo ja sen rakentaminen.
%  Traduction 2026 d'apres texto/io, controlee sur texto/fr et
%  texto/sv. Consigne : voir texto/fi/10-taulukko-01.tex.
%
%  Le tableau 5 est un DIALOGUE, et les deux volumes composent le nom
%  du parleur en petites capitales : on garde \textsc{}.
%
%  LES DEUX NOMS PROPRES SE TRADUISENT, ET L'IDO L'A DEJA FAIT :
%  Leblanc est « Blanko », Roux est « Rufo ». Le finnois forme des
%  noms de famille sur les memes couleurs — valkoinen, ruskea — et
%  ecrit donc herra Valkonen et herra Ruskonen, qui sont de vrais
%  patronymes finnois.
%
%  DEUX INVERSIONS EVITEES D'AVANCE. Au bloc 01-30, l'ido met
%  l'instrument avant l'objet — « per gipso (70) il gipsizas la
%  plafoni (26) » — et le finnois le suit sans peine : l'adessif se
%  porte en tete de phrase comme n'importe quel complement. Au bloc
%  01-49, « la sulo (10) dil extera pordo (8) » preposerait en
%  genitif — « ulko-oven (8) kynnys (10) » — et l'on ecrit donc
%  « kynnyksen (10) yli ulko-oven (8) kohdalla ».
% ===================================================================

%%K t05-apar-1 apar
\VUsaut{6.0mm}
\VUtitre{15.0pt}{60}{TAULUKKO N:o 5}
\VUsaut{1.4mm}
\VUtitre{11.4pt}{0}{Talo ja sen rakentaminen.}
\VUsaut{2.2mm}

%%K t05-tit-1 sub
\VUpk{10.2pt}{40}{Hyvä kohtaaminen.}
\VUsaut{3.0mm}

%%K t05-01-1 p
\textsc{Juhani}. --- Setä, haluaisin kovasti tietää, kuinka \VUgras{talo}
rakennetaan.

%%K t05-01-2 p
\textsc{Setä} (80). --- Se on hyvin helppoa, ystäväni, tule mukaani, näytän
sinulle sen, jota herra Bernhard (79) rakennuttaa ja jossa minä tulen asumaan.

%%K t05-01-3 p
\textsc{Juhani}. --- Kuka sitten on piirtänyt tämän talon
\VUgras{piirustuksen} \textsuperscript{(76)}?

%%K t05-01-4 p
\textsc{Setä}. --- \VUgras{Arkkitehti} \textsuperscript{(75)}; hän esitti hyvin
selvän piirroksen, joka hyväksyttiin, ja \VUgras{rakennusmestari}
\textsuperscript{(77)} aloitti työt.

%%K t05-01-5 p
\textsc{Juhani}. --- Kuka sitten on rakennusmestari?

%%K t05-01-6 p
\textsc{Setä}. --- Hän on herra Valkonen, ystäväsi Jaakobin isä. Hän johtaa
työmiehiä yhdessä herra Ruskosen, arkkitehdin, kanssa.

%%K t05-01-7 p
Katso! juuri oikeaan aikaan tulee herra Valkonen \VUgras{taittomittoineen}
\textsuperscript{(78)}, ja herra Ruskonen piirustuksineen.

%%K t05-01-8 p
Hyvää päivää, hyvät herrat. Kuinka voitte?

%%K t05-01-9 p
\textsc{Herra Ruskonen}. --- Hyvin hyvin, kiitos. Hyvää päivää, Juhani, kuinka
tuo lapsi onkaan kasvanut!

%%K t05-01-10 p
\textsc{Setä}. --- Niin, todellakin, hän on pieni mies. Hän on hyvin vakava,
hänelle täytyy selittää kaikki mitä hän näkee. Tänään hän haluaisi tietää,
kuinka talo rakennetaan.

%%K t05-tit-2 sub
\VUpk{10.2pt}{40}{Työmiehet.}
\VUsaut{3.0mm}

%%K t05-01-11 p
\textsc{Herra Valkonen}. --- No niin, tule mukaani, pikku ystävä, selitän sen
sinulle heti, sillä pidän kovasti lapsista, jotka tahtovat oppia.

%%K t05-01-12 p
Näet tuon työmiehen, joka pitää \VUgras{lapiota} \textsuperscript{(82)}
kädessään : hän on \VUgras{apumies} \textsuperscript{(81)}. Hän kaivaa
\VUgras{ojaa} \textsuperscript{(46)} pannakseen sinne vesijohdon. Hän on myös
kaivanut maata siltä paikalta, jossa talo on, jotta muurari voisi pystyttää
neljä \VUgras{muuria}. Sitä osaa noista muureista, joka on maassa, sanotaan
\VUgras{perustukseksi} \textsuperscript{(2)}. Tila, joka jää perustusten väliin,
muodostaa \VUgras{kellarin} \textsuperscript{(3)}. Tuossa kellarissa on pieni
ikkuna, jota sanotaan \VUgras{kellarinikkunaksi} \textsuperscript{(5)},
\VUgras{ovi} \textsuperscript{(4)} ja portaat.

%%K t05-01-13 p
\VUgras{Muurari} \textsuperscript{(108)} jatkoi muurien rakentamista ja jätti
aukot \VUgras{oville} \textsuperscript{(8)} ja \VUgras{ikkunoille}
\textsuperscript{(17)}.

%%K t05-01-14 p
\textsc{Juhani}. --- Mutta tuota muuraria en näe!

%%K t05-01-15 p
Herra \textsc{Valkonen}. --- Nyt hän on lakannut työskentelemästä talossa. Hän
on \VUgras{tallin} \textsuperscript{(54)} luona, \VUgras{telineellään}
\textsuperscript{(110)}, hän rakentaa \VUgras{pesutuvan}
\textsuperscript{(56)} muuria.

%%K t05-01-16 p
Hänellä on muurauslasta, laastiastia ja \VUgras{luotilanka}
\textsuperscript{(109)}. Mutta noita nimiä et muista.

%%K t05-01-17 p
\textsc{Juhani}. --- Kylläpä muistan, sillä pyydän heti sedältäni pientä
muurauslastaa ja laastiastiaa, ja teen itse luotilangan naruun.

%%K t05-tit-3 sub
\VUsaut{3.0mm}
\VUpk{10.2pt}{40}{Rakennusaineet.}
\VUsaut{3.0mm}

%%K t05-01-18 p
\textsc{Setä}. --- Ha! hyvin hyvä. Mutta sanopa minulle, Juhani, mitä tarvitaan
talon rakentamiseen?

%%K t05-01-19 p
\textsc{Juhani}. --- Tarvitaan \VUgras{hakattua kiveä} \textsuperscript{(59)} ja
\VUgras{laastia} \textsuperscript{(65)}.

%%K t05-01-20 p
\textsc{Setä}. --- Aivan niin, katso tuota miestä, jolla on suuri hattu,
\VUgras{kärryillään} \textsuperscript{(136)}. Hän on \VUgras{ajuri}
\textsuperscript{(135)}. Joka päivä hän tuo tänne \VUgras{kivilohkareita}
\textsuperscript{(60)}. Hän purkaa kuormansa pihalla, ja \VUgras{kivenhakkaajat}
\textsuperscript{(83)} hakkaavat lohkareet ja sahaavat ne
\VUgras{kivisahallaan} \textsuperscript{(85)}. \VUgras{Kantaja}
\textsuperscript{(146)} vie ne muurarille, joka asettaa ne paikoilleen.

%%K t05-01-21 p
Sillä aikaa \VUgras{laastinsekoittaja} \textsuperscript{(87)} valmistaa laastin.
Hän tarvitsee \VUgras{hiekkaa} \textsuperscript{(62)}, \VUgras{kalkkia}
\textsuperscript{(63)} ja \VUgras{vettä} \textsuperscript{(64)}. Kalkki on hänen
vieressään \VUgras{säkissä} \textsuperscript{(71)}, \VUgras{muurarin apulainen}
\textsuperscript{(111)} tuo hänelle hiekan \VUgras{kottikärryillä}
\textsuperscript{(112)} ja \VUgras{oppipoika} \textsuperscript{(113)} on pumpun
(58) luona. Hän täyttää \VUgras{sangon} \textsuperscript{(114)}. Laasti on pian
valmis. \VUgras{Laastinkantaja} \textsuperscript{(107)} ottaa sen ja kantaa sen
ylös, tikkaita pitkin, telineelle, jossa työskentelee muurari, joka viimeistelee
talon sitä sivua.

%%K t05-01-22 p
Ja nyt, miksi sanotaan sitä työmiestä, joka tekee \VUgras{katon}
\textsuperscript{(31)}?

%%K t05-01-23 p
\textsc{Juhani}. --- Hän on \VUgras{katonlaskija} \textsuperscript{(118)}.

%%K t05-tit-4 sub
\VUsaut{3.0mm}
\VUpk{10.2pt}{40}{Talon katto.}
\VUsaut{3.0mm}

%%K t05-01-24 p
Herra \textsc{Valkonen}. --- Niin, mutta ensin tulee \VUgras{kirvesmies}
\textsuperscript{(115)}.

%%K t05-01-25 p
Kirvesmies tekee \VUgras{kattotuolin} \textsuperscript{(35)}. Hän liittää
\VUgras{palkit} \textsuperscript{(37)} yhteen \VUgras{puutapeilla}
\textsuperscript{(117)}. Nuo työmiehet, tuolla ylhäällä, väljine
samettihousuineen, ovat kirvesmiehiä. Toinen pitää pientä palkkia, ja toinen
sahaa tapin poikki \VUgras{kirveellään} \textsuperscript{(116)}. Se, joka on
\VUgras{savupiipun} \textsuperscript{(41)} luona, on katonlaskija. Hän naulaa
ruoteet (*) palkkeihin, ja \VUgras{liuskekiven} \textsuperscript{(66)}
ruoteisiin. Joskus hän panee kattotiiliä.

%%K t05-noto-1 noto
\VUnotes{143.2mm}{%
(*) \VUgras{Balk-o} = S. \textit{Balken}, E. \textit{joist}, R. \textit{solive},
I. \textit{travicello}, V. \textit{balka}, E. Port. \textit{viga}. Tekninen
juuri, jota ei tähän asti ole hyväksytty, mutta jota on syytä ehdottaa
kääntämään sen R. sanan \textit{solive} merkitys, joka ei ole palkki R.
\textit{poutre} eikä pieni palkki. Se on nelikulmainen puutanko, joka kannattaa
lattiaa ja kattoa.%
}

%%K t05-01-26 p
Tallin katto on \VUgras{tiilikatteinen} \textsuperscript{(55)}, talon on
\VUgras{liuskekatteinen} \textsuperscript{(32)} ja kuistin on
\VUgras{sinkkiä} \textsuperscript{(24)}.

%%K t05-01-27 p
\textsc{Juhani}. --- Tekeekö katonlaskija myös sinkkikatot?

%%K t05-01-28 p
Herra \textsc{Valkonen}. --- Ei, vaan \VUgras{sinkkiseppä}
\textsuperscript{(121)} eli \VUgras{lyijyseppä}, se joka panee
\VUgras{kattoikkunan} \textsuperscript{(38)} talon katolle. Hän panee myös
\VUgras{syöksytorvet} \textsuperscript{(43)} ja \VUgras{räystäskourut}
\textsuperscript{(42)}. Hän juottaa sinkin ja pellin. Hän on kiinnittänyt
\VUgras{ukkosenjohtimen} \textsuperscript{(40)} ja \VUgras{tuuliviirin}
\textsuperscript{(39)} \VUgras{päätykolmioon} \textsuperscript{(36)}.

%%K t05-01-29 p
\textsc{Juhani}. --- Onko talo sitten rakennettu valmiiksi?

%%K t05-tit-5 sub
\VUsaut{3.0mm}
\VUpk{10.2pt}{40}{Työkalut.}
\VUsaut{3.0mm}

%%K t05-01-30 p
Herra \textsc{Valkonen}. --- Ei aivan. \VUgras{Rappari}
\textsuperscript{(99)} rakentaa \VUgras{tiilistä} \textsuperscript{(61)} ne
väliseinät, jotka jakavat sen eri huoneisiin. \VUgras{Kipsillä}
\textsuperscript{(70)} hän rappaa \VUgras{katot} \textsuperscript{(26)} ja
muurit. Nyt hän työskentelee \VUgras{kuistin} \textsuperscript{(33)} katolla.
Hänellä on, kuten muurarilla, \VUgras{muurauslasta} \textsuperscript{(100)} ja
\VUgras{laastiastia} \textsuperscript{(101)}.

%%K t05-01-31 p
Sitten puuseppä tekee ovet, ikkunat, \VUgras{parkettilattiat}
\textsuperscript{(16)} ja \VUgras{ikkunaluukut} \textsuperscript{(19)}.

%%K t05-01-32 p
Nyt hän työskentelee pihalla \VUgras{höyläpenkkinsä} \textsuperscript{(90)}
ääressä. Hänellä on \VUgras{saha} \textsuperscript{(94)} puun (69) sahaamiseen,
\VUgras{höylä} \textsuperscript{(91)}, \VUgras{ruuvipenkki}
\textsuperscript{(92)}, \VUgras{talttarauta} \textsuperscript{(94 bis)},
\VUgras{harppi} \textsuperscript{(95)} ja puinen \VUgras{nuija}
\textsuperscript{(95 bis)}.

%%K t05-01-33 p
\textsc{Juhani}. --- Ha! mutta puusepäntyö on hauskempaa kuin muuraaminen. Setä,
sinun täytyy ostaa minulle höyläpenkki, saha ja höylä, sillä teen pian kopin
Medorille.

%%K t05-01-34 p
\textsc{Setä}. --- Kyllä, poikaseni.

%%K t05-01-35 p
\textsc{Juhani}. --- Kuinka iloinen olenkaan! Ja kuka on se, joka on ulko-oven
luona?

%%K t05-tit-6 sub
\VUsaut{3.0mm}
\VUpk{10.2pt}{40}{Maalaus.}
\VUsaut{3.0mm}

%%K t05-01-36 p
Herra \textsc{Valkonen}. --- Hän on \VUgras{maalari} \textsuperscript{(102)}.
Hän on tikkaillaan \VUgras{maalipurkki} \textsuperscript{(103)} ja
\VUgras{sivellin} \textsuperscript{(104)} kädessään. Hän maalaa ulko-oven puuta
säilyttääkseen sen kauemmin.

%%K t05-01-37 p
Hän myös liimaa tapetit kerroksiin.

%%K t05-01-38 p
\VUgras{Tapetoija} \textsuperscript{(107)}, joka on \VUgras{tikkailla}
\textsuperscript{(106)}, \VUgras{parvekkeella} \textsuperscript{(28)},
kiinnittää \VUgras{markiisin} \textsuperscript{(29)}.

%%K t05-01-39 p
Työmies, joka on suuren ikkunan luona, on \VUgras{lasittaja}
\textsuperscript{(96)}. Hän kiinnittää \VUgras{ikkunaruudut}
\textsuperscript{(18)} \VUgras{kitillä} \textsuperscript{(73)}. Toinen työmies,
joka polvistuu kellarin oven luona, on \VUgras{lukkoseppä}
\textsuperscript{(97)}. Hän asettaa paikoilleen \VUgras{lukon}
\textsuperscript{(6)}. Hänen \VUgras{viilansa} \textsuperscript{(98)} on hänen
vieressään.

%%K t05-01-40 p
\textsc{Juhani}. --- Onko se työmies, joka lyö \VUgras{vasarallaan}
\textsuperscript{(84)} rautakappaletta, myös lukkoseppä?

%%K t05-01-41 p
Herra \textsc{Valkonen}. --- Oikeammin sanoen, hän on seppä (125). Hän takoo
\VUgras{rautatyöt} \textsuperscript{(67)} \VUgras{sähköasentajalle}
\textsuperscript{(124)}, joka on \VUgras{vaunuvajan} \textsuperscript{(51)}
katolla. Hänellä on \VUgras{ahjo} \textsuperscript{(126)}, \VUgras{alasin}
\textsuperscript{(127)} ja \VUgras{pihdit} \textsuperscript{(128)}.

%%K t05-01-42 p
Sähköasentaja kiinnittää puhelimen \VUgras{kuparisen} \textsuperscript{(44)}
\VUgras{langan}.

%%K t05-tit-7 sub
\VUpk{10.2pt}{40}{Puutarhanhoito.}
\VUsaut{3.0mm}

%%K t05-01-43 p
\textsc{Setä}. --- \VUgras{Pihan} \textsuperscript{(45)} perällä,
\VUgras{ristikkoaidan} \textsuperscript{(47)} ja \VUgras{portin}
\textsuperscript{(48)} luona näet \VUgras{puutarhurin} \textsuperscript{(129)}.
Hän laittaa kuntoon \VUgras{pientä puutarhaa} \textsuperscript{(49)}. Hän on
kääntänyt maan \VUgras{lapiollaan} \textsuperscript{(131)}, hän kylvää siemeniä
ja haravoi \VUgras{haravalla} \textsuperscript{(130)}. Sitten hän
\VUgras{kastelukannullaan} \textsuperscript{(132)} kastelee
\VUgras{kukkapenkit} \textsuperscript{(150)} ja taimet kasvavat.

%%K t05-01-44 p
\VUgras{Tallirenki} \textsuperscript{(133)}, joka on juuri hoitanut hevosen ja
antanut sille rehua, puhdistaa \VUgras{vaunut} \textsuperscript{(52)} sienellä,
\VUgras{käsipumpulla} \textsuperscript{(134)} ja sangollisella vettä. Sitten hän
ajaa ne takaisin vaunuvajaan ja sulkee oven paksulla \VUgras{salvalla}
\textsuperscript{(53)}.

%%K t05-01-45 p
Kas niin, olet tyytyväinen, luulen ma, olet saanut riittävät selitykset.

%%K t05-01-46 p
\textsc{Juhani}. --- Kyllä, setä, mutta näkisin mielelläni talon sisäpuolen.

%%K t05-01-47 p
\textsc{Setä}. --- Se tapahtuu huomenna. Herra Ruskonen selittää sinulle
piirustuksen ja mainitsee kunkin huoneen nimen.

%%K t05-tit-8 sub
\VUsaut{3.0mm}
\VUpk{10.2pt}{40}{Talon sisäinen jako.}
\VUsaut{2.4mm}

%%K t05-apar-2 apar
{\centering\textit{(Katso piirustusta.)}\par}
\VUsaut{2.4mm}

%%K t05-01-48 p
\textsc{Arkkitehti}. --- Tule, Juhani, annan sinun heti tutkia talon eri osat.

%%K t05-01-49 p
Menkäämme sisään \VUgras{pohjakerrokseen} \textsuperscript{(7)}. Tässä on
\VUgras{soittokellon} \textsuperscript{(11)} nappi. Nousemme kolme
\VUgras{askelmaa} \textsuperscript{(14)} \VUgras{ulkoportailla}
\textsuperscript{(9)}, kuljemme \VUgras{kynnyksen} \textsuperscript{(10)} yli
\VUgras{ulko-oven} \textsuperscript{(8)} kohdalla, ja tässä seisomme
\VUgras{portaiden} \textsuperscript{(13)} juurella.

%%K t05-01-50 p
\textsc{Juhani}. --- Se on käytävä.

%%K t05-01-51 p
\textsc{Arkkitehti}. --- Ei, se on \VUgras{eteishalli} \textsuperscript{(12)}.
\VUgras{Käytävä} \textsuperscript{(a)} on kauempana. Oikealla, tässä on
\VUgras{sali} \textsuperscript{(b)} ja \VUgras{ruokasali}
\textsuperscript{(c)}, ruokasalia vastapäätä on \VUgras{keittiö}
\textsuperscript{(d)} ja \VUgras{tarjoiluhuone} \textsuperscript{(e)}, sitten
edimpänä \VUgras{työhuone} \textsuperscript{(f)}. \VUgras{Kuisti}
\textsuperscript{(23)} \VUgras{lasiovineen} \textsuperscript{(25)} on salin
vieressä.

%%K t05-01-52 p
Nousemme heti portaita. Pidä kiinni \VUgras{kaiteesta}
\textsuperscript{(15)} ja laske askelmat. Niitä on neljätoista. Tässä on
\VUgras{porrastasanne} \textsuperscript{(m)}, olemme ensimmäisessä
\VUgras{kerroksessa} \textsuperscript{(27)}.

%%K t05-tit-9 sub
\VUsaut{3.0mm}
\VUpk{10.2pt}{40}{Ensimmäinen kerros.}
\VUsaut{3.0mm}

%%K t05-01-53 p
Portaita vastapäätä ovat \VUgras{käymälä} \textsuperscript{(20)} ja
\VUgras{pukuhuone} \textsuperscript{(j)}. Oikealla kaunis \VUgras{makuuhuone}
\textsuperscript{(g)}, jossa on kaksi ikkunaa talon \VUgras{julkisivulle}
\textsuperscript{(1)}. Vasemmalla ovat \VUgras{kylpyhuone}
\textsuperscript{(1)} ja \VUgras{vierashuone} \textsuperscript{(i)}, jossa on
suuri \VUgras{alkovi} \textsuperscript{(h)} : Makuuhuoneen vieressä on
\VUgras{lastenhuone} \textsuperscript{(k)}. Se on tilavan \VUgras{terassin}
\textsuperscript{(21)} vieressä. Tuon terassin alla on toinen käymälä (20) ja,
muurilla, tässä on \VUgras{kello} \textsuperscript{(22)}, joka soittaa
päivälliselle.

%%K t05-01-54 p
\textsc{Juhani}. --- Nouskaamme \VUgras{toiseen kerrokseen}.

%%K t05-01-55 p
\textsc{Arkkitehti}. --- Toista kerrosta ei ole. Katon alla on vain
\VUgras{ullakkohuoneita} \textsuperscript{(33)} ja vintti (34). Olemme
lopettaneet tutkimuksen, voimme mennä taas alas.

%%K t05-01-56 p
\textsc{Juhani}. --- Kiitos, herra, se on hyvin mielenkiintoista. Kerron heti
isälleni kaiken mitä olen nähnyt.
\VUsaut{4.0mm}
\VUfilet{22mm}
